\section{Related Work}


\paragraph{Memory-Augmented Large Language Models}
A growing body of work focuses on building latent-space memory systems, as characterized in M+\citep{mplus}, where transformer architectures are modified to support memory augmentation. These memory components can reside in various latent forms, including model parameters\citep{self-param}, external memory matrices~\citep{larimar,infini-attention}, hidden states~\citep{RMT, he2024camelot, memoryllm, mplus}, soft prompts~\citep{memoryTransformer}, and key-value caches~\citep{MemoryBank,cartridges,snapkv,h2o}. While these approaches demonstrate promising results and continue to advance the frontier of latent-space memory, most require retraining the model~\citep{memoryllm, mplus, infini-attention}, making them incompatible with powerful closed-source models like GPT-4 or DeepSeek-R1. Additionally, methods based on key-value caching rely on preserving past keys and values, functioning more as long-context methods rather than true memory systems that support abstraction, consolidation, and reasoning over stored experiences~\citep{lscs}.


\paragraph{Memory-Augmented LLM Agents}
Token-level memory remains the dominant approach in current LLM agents~\citep{mplus}, where past conversational content is stored in raw text form within external databases. Notable examples include commercial systems such as Zep~\citep{zep}, Mem0~\citep{mem0}, and MemGPT~\citep{memgpt}. These agents perform well on long-term conversational benchmarks~\citep{locomo,longmemeval} and document-based retrieval tasks. However, they often fall short in real-world applications due to simplistic memory architectures. Most notably, the absence of modular memory components hinders effective memory routing and leads to inefficiencies in retrieval and usage.


\paragraph{Various Memory Types}
Cognitive science broadly categorizes memory into short-term (working) memory and long-term memory~\citep{sridhar2023cognitive}. In the context of LLMs, short-term memory is often mapped to the input context window~\citep{li2024memory}, while long-term memory becomes a catch-all category for any information outside the context~\citep{mplus}. To address this limitation, recent works have proposed finer-grained memory architectures. For example, \citet{episodic_memory_missing_piece} emphasizes the importance of episodic memory in LLM agents. Other systems incorporate both episodic and semantic memory~\citep{arigraph,memory_consciousness_and_llm,tulving1985memory}. Semantic memory has also been highlighted as critical for real-world reasoning and abstraction~\citep{procedural_memory_all_you_need,kim2023machine}. In addition, procedural memory, responsible for learned skills and routine tasks, has been identified as another crucial component~\citep{procedural_memory_all_you_need}. Despite these advancements, existing methods stop at identifying individual memory types, and they are not formed into a comprehensive memory system.


\paragraph{Multi-Agent Systems}
Rather than relying on a monolithic agent, recent advances explore multi-agent frameworks where specialized agents coordinate to accomplish complex tasks. Early systems like AutoGPT and BabyAGI~\citep{AutoGPT,BabyAGI} adopt an autonomous planning-execution loop while maintaining a shared memory log. More recent designs introduce role specialization: MetaGPT~\citep{Hong2023} mimics a software development team structure, and AgentVerse~\citep{Chen2024} assigns agents to specific roles such as planning or evaluation. Cognitive theories also support modularity in memory, particularly distinctions between episodic and semantic types, as emphasized in \citet{Liao2024}. \ours builds on these ideas by deploying eight specialized agents, each managing a distinct memory type (e.g., episodic, semantic, procedural), and coordinating to process multi-modal inputs effectively.
